
\documentclass{article}
\usepackage{colortbl}
\usepackage{makecell}
\usepackage{multirow}
\usepackage{supertabular}

\begin{document}

\newcounter{utterance}

\centering \large Interaction Transcript for game `mmlu\_pro', experiment `mmlu\_pro\_health', episode 682 with reflexion{-}Qwen3.5{-}27B.
\vspace{24pt}

{ \footnotesize  \setcounter{utterance}{1}
\setlength{\tabcolsep}{0pt}
\begin{supertabular}{c@{$\;$}|p{.15\linewidth}@{}p{.15\linewidth}p{.15\linewidth}p{.15\linewidth}p{.15\linewidth}p{.15\linewidth}}
    \# & \multicolumn{2}{c}{Player} && \multicolumn{2}{c}{Game Master} \\
    \hline

    \theutterance \stepcounter{utterance}  
    & & & \multicolumn{4}{p{0.6\linewidth}}{
        \cellcolor[rgb]{0.9,0.9,0.9}{
            \makecell[{{p{\linewidth}}}]{
                \texttt{\tiny{[P1$\langle$GM]}}
                \texttt{You have attempted this same task before and did not succeed. These are the plans you wrote after those attempts.} \\
\texttt{Trial \#1: New plan: I incorrectly identified the diagnosis as Munchausen syndrome (E) by focusing solely on the history of hospitalizations and dramatic symptoms, while overlooking the critical detail that "definitive causes... have not been found despite extensive diagnostic studies." Munchausen syndrome (Factitious Disorder Imposed on Self) requires the intentional fabrication or induction of symptoms, which is not explicitly stated here. The description of long{-}standing, unexplained physical symptoms with extensive medical workups and dramatic presentation is the classic definition of Somatic Symptom Disorder (A). In the future, I will carefully distinguish between disorders involving intentional deception (Factitious/Munchausen) and those involving genuine distress over unexplained physical symptoms (Somatic Symptom Disorder) by checking for explicit evidence of fabrication versus a history of negative diagnostic findings. I will select option A for this specific clinical vignette.} \\
\\ 
\texttt{Trial \#2: New plan: I previously selected Munchausen syndrome (E) by overemphasizing the dramatic presentation and hospitalization history, while failing to distinguish between intentional fabrication and genuine distress over unexplained symptoms. The key differentiator is the absence of evidence for intentional deception; the prompt states "definitive causes... have not been found," which aligns with Somatic Symptom Disorder (A), where the distress is real but the cause is unknown, rather than Factitious Disorder where symptoms are feigned. In future attempts, I will explicitly check for keywords indicating intentional lying or self{-}harm (e.g., "faking," "inducing," "lying") before selecting Munchausen. If the vignette describes a long history of unexplained physical symptoms and extensive medical seeking without evidence of fabrication, I will prioritize Somatic Symptom Disorder (A) over Munchausen syndrome.} \\
\\ 
\texttt{Trial \#3: New plan: I previously selected Munchausen syndrome (E) by misinterpreting the "dramatic and exaggerated descriptions" and "conflicting statements" as definitive proof of intentional fabrication. However, the prompt explicitly states that "definitive causes... have not been found despite extensive diagnostic studies," which is the hallmark of Somatic Symptom Disorder (A), where the patient genuinely experiences distress from unexplained symptoms. Munchausen syndrome (Factitious Disorder) requires evidence of *intentional* production or feigning of symptoms, which is not confirmed here; the conflicting statements are likely due to the patient's distress and memory issues rather than calculated deception. In future attempts, I will strictly differentiate between "unexplained symptoms with negative workups" (Somatic Symptom Disorder) and "intentional fabrication" (Munchausen). If the vignette lacks explicit evidence of the patient lying or inducing illness, I will default to Somatic Symptom Disorder (A) for cases involving chronic, unexplained physical complaints and excessive medical seeking.} \\
\\ 
\\ 
\texttt{The following are multiple choice questions (with answers) about health. Answer only with "The answer is (X)" where X is the correct letter choice.} \\
\\ 
\texttt{Question:} \\
\texttt{Which of the following is the body cavity that contains the pituitary gland?} \\
\texttt{Options:} \\
\texttt{A. Ventral} \\
\texttt{B. Dorsal} \\
\texttt{C. Buccal} \\
\texttt{D. Thoracic} \\
\texttt{E. Pericardial} \\
\texttt{F. Abdominal} \\
\texttt{G. Spinal} \\
\texttt{H. Pelvic} \\
\texttt{I. Pleural} \\
\texttt{J. Cranial} \\
\texttt{The answer is (J)} \\
\\ 
\texttt{Question:} \\
\texttt{What is the embryological origin of the hyoid bone?} \\
\texttt{Options:} \\
\texttt{A. The third and fourth pharyngeal arches} \\
\texttt{B. The fourth pharyngeal arch} \\
\texttt{C. The third pharyngeal arch} \\
\texttt{D. The second pharyngeal arch} \\
\texttt{E. The second, third and fourth pharyngeal arches} \\
\texttt{F. The first pharyngeal arch} \\
\texttt{G. The second and third pharyngeal arches} \\
\texttt{H. The first and third pharyngeal arches} \\
\texttt{I. The first, second and third pharyngeal arches} \\
\texttt{J. The first and second pharyngeal arches} \\
\texttt{The answer is (G)} \\
\\ 
\texttt{Question:} \\
\texttt{What is the difference between a male and a female catheter?} \\
\texttt{Options:} \\
\texttt{A. Female catheters are used more frequently than male catheters.} \\
\texttt{B. Male catheters are bigger than female catheters.} \\
\texttt{C. Male catheters are more flexible than female catheters.} \\
\texttt{D. Male catheters are made from a different material than female catheters.} \\
\texttt{E. Female catheters are longer than male catheters.} \\
\texttt{F. Male catheters are longer than female catheters.} \\
\texttt{G. Female catheters are bigger than male catheters.} \\
\texttt{H. Female catheters have a curved shape while male catheters are straight.} \\
\texttt{I. Male and female catheters are different colours.} \\
\texttt{J. Male catheters have a smaller diameter than female catheters.} \\
\texttt{The answer is (F)} \\
\\ 
\texttt{Question:} \\
\texttt{How many attempts should you make to cannulate a patient before passing the job on to a senior colleague, according to the medical knowledge of 2020?} \\
\texttt{Options:} \\
\texttt{A. 1} \\
\texttt{B. Unlimited attempts} \\
\texttt{C. 5} \\
\texttt{D. 0, always pass on to a senior colleague} \\
\texttt{E. 7} \\
\texttt{F. 2} \\
\texttt{G. 4} \\
\texttt{H. 6} \\
\texttt{I. 3} \\
\texttt{J. 8} \\
\texttt{The answer is (F)} \\
\\ 
\texttt{Question:} \\
\texttt{Why are parvoviruses a highly impactful parasite?} \\
\texttt{Options:} \\
\texttt{A. They are able to alter the host's DNA} \\
\texttt{B. Because they have no nucleic acid} \\
\texttt{C. They can survive in extreme temperatures} \\
\texttt{D. Only replicate in dividing cells} \\
\texttt{E. They can infect multiple species} \\
\texttt{F. They don't require a host to survive} \\
\texttt{G. Can integrate into host chromosomes} \\
\texttt{H. N/A} \\
\texttt{I. N/A} \\
\texttt{J. N/A} \\
\texttt{The answer is (B)} \\
\\ 
\texttt{Question:} \\
\texttt{A 50{-}year{-}old woman comes to the office for the first time because of recurrent abdominal pain. Review of her extensive medical chart, which she has brought with her, discloses that she has a long history of varying physical symptoms. Definitive causes for these symptoms have not been found despite extensive diagnostic studies, consultations with many physicians, and several surgical explorations. She gives dramatic and exaggerated descriptions of her present and past symptoms, and she makes conflicting statements about her history. She has been hospitalized at least 23 times since age 18 years. Which of the following is the most likely diagnosis?} \\
\texttt{Options:} \\
\texttt{A. Somatic symptom disorder} \\
\texttt{B. Generalized anxiety disorder} \\
\texttt{C. Dependent personality disorder} \\
\texttt{D. Borderline personality disorder} \\
\texttt{E. Munchausen syndrome} \\
\texttt{F. Conversion disorder} \\
\texttt{G. Schizophrenia} \\
\texttt{H. Occult medical disorder} \\
\texttt{I. Narcissistic personality disorder} \\
\texttt{J. Histrionic personality disorder} \\
            }
        }
    }
     \\ \\

    \theutterance \stepcounter{utterance}  
    & \multicolumn{4}{p{0.6\linewidth}}{
        \cellcolor[rgb]{0.561,0.737,0.561}{
            \makecell[{{p{\linewidth}}}]{
                \texttt{\tiny{[P1$\rangle$GM]}}
                \texttt{The answer is (A)} \\
            }
        }
    }
    & & \\ \\

    \theutterance \stepcounter{utterance}  
    & & & \multicolumn{2}{p{0.3\linewidth}}{
        \cellcolor[rgb]{0.9,0.9,0.9}{
            \makecell[{{p{\linewidth}}}]{
                \texttt{\tiny{[GM$|$GM]}}
                \texttt{A} \\
            }
        }
    }
    & & \\ \\

    \theutterance \stepcounter{utterance}  
    & & & \multicolumn{2}{p{0.3\linewidth}}{
        \cellcolor[rgb]{0.9,0.9,0.9}{
            \makecell[{{p{\linewidth}}}]{
                \texttt{\tiny{[GM$|$GM]}}
                \texttt{A} \\
            }
        }
    }
    & & \\ \\

    \theutterance \stepcounter{utterance}  
    & & & \multicolumn{2}{p{0.3\linewidth}}{
        \cellcolor[rgb]{0.9,0.9,0.9}{
            \makecell[{{p{\linewidth}}}]{
                \texttt{\tiny{[GM$|$GM]}}
                \texttt{game\_result = WIN} \\
            }
        }
    }
    & & \\ \\

\end{supertabular}
}

\end{document}
